\appendix
\section{Running and extending the bundle}\label{sec:running}

The standard-library-only replay command is:
\begin{lstlisting}
cd prototype
python -S replay_all.py ../results/run1
\end{lstlisting}
It checks the 90 stored construction certificates against separately stored expected queries, then checks the 643 concrete witness traces against their verified source machines. Replaying the corpus does not regenerate it or import the producer.

To run all experiments into a fresh directory and run the small interface suite:
\begin{lstlisting}
python -S experiments.py --out ../reproduced-results
python -S -m unittest test_interfaces
\end{lstlisting}
Use a new output directory to preserve historical measurements. Reusing an existing directory overwrites its generated files; the research script does not enforce a no-overwrite policy. Do not run the tests with Python's \code{-O} flag, which removes assertions used by the experiment harness. The independent checker uses explicit rejection checks rather than relying on assertions.

To regenerate and check one certificate:
\begin{lstlisting}
python -S producer.py \
  ../results/run1/queries/successor_cut.json \
  ../reproduced-successor-cut.json
python -S checker.py \
  ../results/run1/queries/successor_cut.json \
  ../reproduced-successor-cut.json
\end{lstlisting}
The producer's \code{--max-states} flag can exercise a budget refusal; \code{--no-minimize} disables quotient generation. The checker expects the query as a separate argument. A bundle claiming \code{unknown} is deliberately not accepted as a proof certificate.

A minimal Python use of least-witness synthesis is:
\begin{lstlisting}[language=Python]
from formula import And, le, pow2, least_graph
from producer import Compiler, extract_witness
from checker import verify, check_witness

relation = And(pow2("y"), le({"x": 1, "y": -1}, -1))
graph = least_graph(relation, "y", "z")
query = {"context": ["x", "y"], "formula": graph}
certificate = Compiler().bundle(graph, query["context"])
verify(certificate, query)
machine = certificate["nodes"][certificate["root"]]["dfa"]
inputs = [10**100 + 37]
record = extract_witness(machine, inputs)
if record is None:
    raise RuntimeError("expected a witness for this total relation")
check_witness(machine, inputs, record)
print(record["witness"])
\end{lstlisting}
The \code{verify} result for this graph need not say its universal closure is valid; nearly all pairs are outside a function graph. What witness extraction needs is the checked graph semantics, not universality of the graph over all input/output pairs.

\section{Certificate schema at a glance}
The complete machine-readable examples are the stored JSON corpus. The following is an outline, not JSON to submit directly:
\begin{lstlisting}
{
  "schema": "forge-auto-1",
  "query": {"context": [...], "formula": ...},
  "nodes": [
    {"kind": "atom", "ctx": [...], "formula": ...,
     "dfa": {"k": ..., "start": ..., "trans": [...],
             "final": [...]},
     "labels": [...]},
    ...
  ],
  "root": <node index>,
  "verdict": {"kind": "valid", "invariant": [...]}
}
\end{lstlisting}
The alternative verdict is \code{counterexample} with \code{word} and \code{values}. A quantified node has \code{child}, \code{subsets}, \code{tail_rank}, and \code{tail_choice}; its formula is an \code{exists} node and its child's context appends exactly that binder. A quotient node has a \code{mapping} indexed by every child state. Products have \code{left}, \code{right}, and \code{pairs}. Negation has a \code{child}.

Every formula atom uses explicit coefficient maps or variable names. Coefficient zero may be retained in a map; it does not alter the mathematics, but the exact supplied query must still match. All numbers on the wire are JSON integers, not floats or decimal strings. This is a small research format; long-term canonical serialization, compatibility versioning, and streaming parsing remain future work.

\section{Review ledger and ownership of claims}
\begin{longtable}{@{}p{0.29\linewidth}p{0.63\linewidth}@{}}
\toprule
Claim or component & Evidence and qualification\\
\midrule\endhead
Forge baseline & Pinned merged status, README, closure chapter and Boolean chapter; not a rerun of the repository's historical experiments.\\
Automata arithmetic & Established prior work; not claimed as a new mathematical method.\\
Correct quantified projection & Mathematical proof in this article; Python implementation and receipts tested; related existing Lean source acknowledged.\\
Canonical witness graph & Mathematical least-choice argument and executed five-family prototype; not an arbitrary higher-order Lean synthesizer.\\
Independent replay & Fresh \code{python -S} process, producer modules not imported; checker not formally verified.\\
Lean library reuse & Documentation and selected source inspection; no local compilation, dependency migration, or axiom audit.\\
Leant/Djex inspiration & Source-bound candidate discipline; no claim to close Church indexing, universe-polymorphism or higher-rank synthesis milestones.\\
Performance & Local single-run Python measurements and a small ablation; no speedup against Lean tactics.\\
Hardening & Several schema checks and mutation controls; no security or denial-of-service certification.\\
\bottomrule
\end{longtable}

The full Forge commit used for the comparison is
\begin{center}\code{58ea206bd7ad501b930add568151217bcc7f82a2}.\end{center}
The companion repositories' branch references observed during review were:
\begin{center}
\begin{tabular}{ll}
Leant & \code{6bf05ad78c467989e68290f2d08bbed40802d485}\\
Djex & \code{e8778f4ebd63e1f9b9b410fa4de8d14a8a04c9e5}\\
Aeacu2 & \code{8b5fb2dd7e4f32e17f6c4dd43f2d8980649ebc42}
\end{tabular}
\end{center}
The Leant and Djex READMEs were retrieved through the main-branch interface; these recorded branch references aid follow-up but are not a proof of a frozen multi-repository build. The inspected Aeacu2 files were retrieved at the listed commit.
